% Rol: Estadísticas.
\chapter{Reportes y estadísticas}\label{cap:estadisticas}

\begin{ficharol}{Estadísticas}{Este capítulo es para usted si prepara cifras e informes
  sobre las admisiones.}
Con este rol usted puede generar cuatro tipos de reporte, filtrarlos, verlos en pantalla y
descargarlos en PDF o CSV (para Excel). Además puede consultar la tabla de carreras. No puede
modificar admisiones.
\end{ficharol}

\section{Los cuatro tipos de reporte}

\begin{tabularx}{\linewidth}{@{}>{\leavevmode\color{ulaazuloscuro}\sffamily\bfseries}p{3.6cm} X@{}}
  \toprule
  \normalfont\bfseries Tipo & \bfseries Para qué sirve \\
  \midrule
  Análisis general (datos) & Una tabla con una fila por admisión y las columnas que usted
                             escoja. Ideal para descargar en CSV y trabajar en Excel. \\
  Consolidado de revisión & Cuántas admisiones hay por día, por estado y por carrera, con
                            totales por día y total general. \\
  Estadísticas & Cuadros resumen: por estado, por día, semana y mes, por carrera, facultad,
                 modalidad, período, género y más. \\
  Reconciliación (embudo) & Compara los asignados contra los procesados: cuántos faltan por
                            atender, cuántos se devolvieron y volvieron, y cuánto tiempo llevan
                            pendientes. \\
  \bottomrule
\end{tabularx}

\section{Armar un reporte}

Entre a \menu{Estadísticas}.

\begin{captura}{estadisticas-form}{2732}{2620}[0.92\linewidth]
  \anota{1}{876}{151}{1314}{226}
  \anota{2}{1410}{151}{1849}{226}
  \anota{3}{650}{322}{1349}{399}
  \anota{4}{1422}{333}{2101}{397}
  \anota{5}{1061}{446}{1664}{526}
  \anota{6}{50}{593}{2677}{1779}
  \anota{7}{50}{1837}{2677}{2273}
  \anota{8}{619}{2320}{2107}{2433}
\end{captura}

\begin{leyenda}
  \item[1] \textbf{Tipo de reporte}: uno de los cuatro de la tabla anterior.
  \item[2] \textbf{Base de fecha}: \emph{Fecha de asignación origen} cuenta a los estudiantes
        asignados en ese rango (vengan o no a inscribirse). \emph{Fecha válida DIGAIPE} cuenta
        lo que se atendió en ese rango.
  \item[3] \textbf{Desde} y \textbf{Hasta}: obligatorias.
  \item[4] \textbf{Rangos rápidos}: llenan las fechas por usted (hoy, últimos 7 días, este mes,
        este año).
  \item[5] \textbf{Incluir cédula y nombre}: déjelo sin marcar salvo que de verdad necesite
        identificar a los estudiantes.
  \item[6] \textbf{Filtros}: proceso, período, modalidad, estado, facultad, carreras y búsqueda
        por texto. Mantenga pulsada la tecla \texttt{Ctrl} para escoger varios valores de una
        lista.
  \item[7] \textbf{Opciones del tipo de reporte}: columnas (análisis general), desgloses
        (estadísticas) o dimensión de la matriz (reconciliación).
  \item[8] \textbf{Botones}: \boton{Vista previa}, \boton{Descargar PDF}, \boton{Descargar CSV},
        \boton{Copiar enlace} y \enlace{Limpiar}.
\end{leyenda}

\begin{pasos}
  \paso{Escoja el tipo}{En \campo{Tipo de reporte} escoja el que necesita.}
  \paso{Ponga las fechas}{Llene \campo{Desde} y \campo{Hasta}, o use un rango rápido.}
  \paso{Filtre (opcional)}{Escoja proceso, carreras u otros filtros. Si no escoge nada, el
        reporte incluye todo.}
  \paso{Vea el resultado}{Pulse \boton{Vista previa}. El resultado aparece debajo del
        formulario.}
  \paso{Descargue}{Pulse \boton{Descargar PDF} para imprimir o \boton{Descargar CSV} para
        abrirlo en Excel.}
\end{pasos}

\begin{importante}[Faltan las fechas]
Si pulsa \boton{Vista previa} sin fechas, aparece el recuadro rojo \emph{«No se puede generar el
reporte — Seleccione las fechas Desde y Hasta.»}
\end{importante}

\section{Leer el resultado}

\begin{captura}{estadisticas-resultado}{2732}{1220}
  \anota{1}{89}{20}{2634}{210}
  \anota{2}{89}{270}{687}{344}
  \anota{3}{49}{497}{2676}{1179}
\end{captura}

\begin{leyenda}
  \item[1] \textbf{Totales}: cuántas admisiones y cuántos estudiantes distintos entran en el
        reporte. Un estudiante con dos carreras cuenta como dos admisiones y un estudiante.
  \item[2] \textbf{Exportar} el mismo resultado en CSV o PDF.
  \item[3] \textbf{Cuadros}: en el tipo \emph{Estadísticas}, primero los indicadores principales
        y luego un cuadro por cada desglose escogido.
\end{leyenda}

Así se ve la primera página del PDF del tipo \emph{Estadísticas}:

\pantalla{reporte-pdf-1}[0.6\linewidth]

\begin{consejo}[Compartir un reporte]
Después de pulsar \boton{Vista previa}, \boton{Copiar enlace} copia la dirección de ese reporte
con todos sus filtros. Pegue el enlace en un correo: quien tenga permiso de estadísticas verá
exactamente el mismo reporte. Si cambia algún filtro, pulse de nuevo \boton{Vista previa} antes
de copiar el enlace.
\end{consejo}

\begin{nota}[Reportes grandes]
Si el reporte tiene más de 5.000 admisiones, el sistema avisa que la descarga puede tardar y
generar un archivo grande. Use filtros para achicarlo.
\end{nota}

\begin{importante}[Datos personales]
Los reportes con cédula y nombre contienen datos personales de los estudiantes. Guárdelos solo
en equipos de la universidad y no los envíe por correos personales ni redes sociales. Cada
consulta queda registrada en la auditoría.
\end{importante}
